\section{Las tres capas del presente}

Se propone una partición ontológica del presente en tres capas jerárquicas de organización de la persistencia. Cada capa se define por su localización, sus marcadores operacionales y su relación con la antisincronización.

\subsection{Capa 1: Presente local intensificado}

\begin{itemize}
    \item \textbf{Localización}: Módulo individual (componente, variable o región).
    \item \textbf{Marcadores}: Hiper-persistencia ($\tau_s$ significativamente superior al promedio reciente) combinada con trapping estructural (LAM o TT por encima de umbrales Config B, o aceleraciones correspondientes).
    \item \textbf{Características ontológicas}: Alta densidad e intensidad local, mayor adhesión a la configuración actual, incremento de la resistencia al cambio \textit{local}.
    \item \textbf{Relación con la antisincronización}: Puede emerger incluso bajo alta antisincronización, pero su influencia sobre la organización global del sistema permanece limitada mientras persista la fragmentación.
\end{itemize}

\begin{conceptbox}[Definición]
El \textit{presente local intensificado} es el estado en el que un módulo sostiene su configuración de persistencia con \textbf{mayor firmeza y adhesión} de lo habitual para él en ese período, medido por la conjunción de hiper-persistencia y trapping estructural.
\end{conceptbox}

\subsection{Capa 2: Coherencia entre presentes locales}

\begin{itemize}
    \item \textbf{Localización}: Relaciones entre módulos.
    \item \textbf{Marcadores}: Reducción significativa de la antisincronización $A(k)$, es decir, alineación de las escalas de persistencia $\tau_s$ de los distintos componentes.
    \item \textbf{Características ontológicas}: Disminución de la autonomía de los presentes locales, reducción de la fragmentación, incremento de la permeabilidad relacional.
    \item \textbf{Rol funcional}: Condición necesaria (aunque no suficiente) para que una intensificación local (Capa 1) pueda propagarse y generar efectos a escala del sistema (Capa 3).
\end{itemize}

La Capa 2 no es un estado de sincronización perfecta ni de identidad de trayectorias, sino de \textit{coherencia de escalas de persistencia}. Los módulos pueden evolucionar de forma distinta, pero lo hacen con ritmos de retención ordinal comparables.

\subsection{Capa 3: Reorganización estructural del sistema}

\begin{itemize}
    \item \textbf{Localización}: El sistema considerado como totalidad.
    \item \textbf{Marcadores}: Cambios detectables en la matriz de coherencias (norma de Frobenius entre matrices consecutivas), aparición de picos hiperestructurales con adelanto significativo respecto a transiciones inyectadas o detectadas, y reorganización global de las relaciones entre módulos.
    \item \textbf{Características ontológicas}: Emergencia de una nueva estructura de relaciones que no se reduce a la suma de los presentes locales intensificados.
    \item \textbf{Condición de emergencia}: Generalmente requiere la conjunción de intensificación local (Capa 1) y reducción previa de fragmentación (Capa 2).
\end{itemize}

\begin{importantbox}
La transición de la Capa 1 a la Capa 3 no es automática. Una intensificación local que ocurre bajo alta antisincronización (alta fragmentación) tiende a permanecer como un «fragmento ontológico»: el sistema puede absorberla sin alterar su organización global gracias a la «resistencia ascendente». Solo cuando la antisincronización desciende —es decir, cuando opera la «resistencia descendente»— se crean las condiciones de permeabilidad necesarias para que la intensificación local adquiera peso estructural.
\end{importantbox}

\subsection{Jerarquía y no reductibilidad}

Las tres capas no son meros niveles descriptivos. Constituyen una jerarquía ontológica en la que cada capa superior presupone condiciones de la inferior pero no se reduce a ellas. La coherencia relacional (Capa 2) presupone la existencia de presentes locales (Capa 1), pero añade una dimensión relacional irreducible. La reorganización estructural (Capa 3) presupone ambas, pero introduce una estructura global que no puede predecirse a partir del mero listado de intensificaciones locales.

\begin{figure}[htbp]
\centering
\fbox{\parbox{0.88\textwidth}{\centering\small
\textbf{Diagrama conceptual de las tres capas del presente}\\[0.3em]
\textbf{Capa 3} (Reorganización estructural) $\leftarrow$ requiere baja antisincronización $\leftarrow$ \textbf{Capa 2} (Coherencia entre presentes locales) $\leftarrow$ \textbf{Capa 1} (Presente local intensificado: hiper-persistencia + trapping RQA)\\[0.3em]
\textit{La antisincronización actúa como regulador de la fragmentación y de la permeabilidad entre capas (ver texto).}
}}
\caption{Esquema de las tres capas del presente y de la dinámica de resistencia. La antisincronización actúa como regulador de la permeabilidad entre las capas.}
\label{fig:capas}
\end{figure}